% sec7_conclusion.tex
\section{Conclusion}
\label{sec:conclusion}

Under a frozen 500-channel evaluation protocol, direct
reinforcement-learning phase policies for a blocked direct-link RIS
scenario provided no measurable benefit over random phase selection,
across every investigated architecture (continuous and codebook
discretization levels) and thirty independent training runs. Deep
unfolding, by contrast, recovers $95\%$ of the Riemannian-AO reference at
$\sim$9.7~ms decision latency, absorbs 2--4-bit quantization losses
gracefully, and degrades coherently under channel error. The
asymmetry between the two epistemologies is not sample-limited (the
200-draw random oracle sits $8$~bps/Hz above both), it is
\emph{structure}-limited.

We offer this not as a verdict against RL writ large, but against an
evaluation culture in which trained-policy comparisons are reported
without frozen test sets, written acceptance rules, or disciplined
selection rules. The three-layer protocol---AO reference,
model-driven learner, model-free challenger---used here is portable to
any emerging-RIS benchmark, and we release the decision log
(DD-001--DD-033), the raw per-channel artifacts, and the entire
figure/statistics pipeline so that any claimed RL-vs-unfolding
ordering can, from now on, be checked against committed files rather
than prose.

\paragraph{Future work.} Three ranked directions follow directly from
the reported data: (i) reward shaping and state-augmented
observations that inject gradient-sign structure into the bandit
observation, giving model-free learners the same structure unfolding
consumes at design time; (ii) sub-millisecond forward-only inference
on hardware targets versus the measured $9.7$~ms, quantifying latency
on real silicon rather than software proxies; and (iii) extension to
multi-RIS and multi-user scaling, where the AO-reference gap is expected
to widen the way it did for the single-RIS hard case here.
